\documentclass{article}
\usepackage[utf8]{inputenc}
\usepackage[T1]{fontenc}
\usepackage{geometry}
\geometry{
    a4paper,
    left=2.5cm,
    right=2.5cm,
    top=2.5cm,
    bottom=2.5cm
}
\usepackage{graphicx}
\usepackage{float}
\usepackage{hyperref}
\usepackage[english]{babel}
\usepackage{amsmath}
\usepackage{amssymb}
\usepackage{array}
\usepackage{booktabs}
\usepackage{caption}
\usepackage{subcaption}
\usepackage{enumitem}
\usepackage{fancyhdr}
\usepackage{titlesec}
\usepackage{url}
\usepackage{textcomp}
\usepackage{xcolor}
\usepackage{listings}
\usepackage{multirow}
\usepackage{longtable}
\usepackage{tabularx}
\usepackage{times}

% Set up page style
\pagestyle{fancy}
\fancyhf{}
\rfoot{\thepage}
\renewcommand{\headrulewidth}{0pt}

% Title formatting
\titleformat{\chapter}[display]{\normalfont\bfseries}{\chaptertitle\thechapter}{0pt}{\Huge}
\titlespacing*{\chapter}{0pt}{50pt}{40pt}

% Document info
\title{Template Document}
\author{Academic Support Agent}
\date{\today}

\begin{document}
\maketitle

\pagenumbering{roman}

% Geef hierna je titel, naam en de datum in,
% de titelpagina wordt automatisch aangemaakt.
\title{LaTeX:\\Een sjabloon voor een rapport}
\author{Jan Janssens\\jan@jansenss.be}
\date{25 november 2004}
\maketitle


% Als je een index met de tabellen of figuren wenst,
% kan je deze hier uitcommentari ren.
%
%

% de samenvatting:
% Content from abstract.tex would be included here


% de hoofdstukken:
% Content from chapter1.tex would be included here

% Content from chapter2.tex would be included here

% Content from chapter3.tex would be included here

% Content from chapter4.tex would be included here



% Content from appendix1.tex would be included here
\end{document}